\begin{table}[t]
\centering
\caption{Summary Statistics: Japanese and Chinese Men in the United States, 1900--1930}
\label{tab:summary}
\begin{threeparttable}
\small
\begin{tabular}{lcccccccc}
\toprule
 & \multicolumn{4}{c}{Japanese} & \multicolumn{4}{c}{Chinese} \\
\cmidrule(lr){2-5} \cmidrule(lr){6-9}
 & 1900 & 1910 & 1920 & 1930 & 1900 & 1910 & 1920 & 1930 \\
\midrule
$N$ & 23,349 & 58,475 & 56,231 & 52,252 & 80,053 & 60,814 & 45,912 & 47,779 \\
Mean age & 26.8 & 30.6 & 36.4 & 39.9 & 40.4 & 41.9 & 41.8 & 37.5 \\
Sex ratio (M/F) & 24.1 & 8.4 & 2.3 & 1.8 & 23.7 & 17.4 & 10.2 & 5.9 \\
Married (\%) & 18.7 & 26.9 & 54.7 & 55.8 & 39.1 & 42.8 & 50.5 & 47.2 \\
Spouse present (\%) & 1.6 & 9.0 & 37.9 & 45.4 & 2.6 & 3.7 & 6.6 & 11.3 \\
Occupational score & 15.3 & 15.3 & 16.0 & 17.0 & 16.7 & 18.7 & 18.5 & 18.7 \\
Literate (\%) & 84.3 & 90.8 & 91.3 & 92.2 & 75.7 & 83.9 & 81.5 & 80.6 \\
Farm owner (\%) & 2.1 & 3.4 & 2.8 & 5.8 & 1.8 & 1.2 & 0.9 & 0.8 \\
\bottomrule
\end{tabular}
\begin{tablenotes}[flushleft]
\footnotesize
\item \textit{Notes:} Full-count U.S.\ Census data from IPUMS, men aged 15--65. ``Married'' includes both spouse present and absent. ``Spouse present'' requires co-residence. Sex ratio is men per woman. OCCSCORE assigns 1950 median income to each occupation. ``Farm owner'' denotes residing on a farm in an owned dwelling.
\end{tablenotes}
\end{threeparttable}
\end{table}
